\section{Contributions}\label{sec:introduction:contributions}
This thesis makes the first end-to-end formal demonstration that a quantum-implicit programming language is achievable.
The language is formally specified, proven type-safe and compiled to execute on gate-based quantum hardware.
The four incorporated publications constitute an integrated body of work; no single publication delivers this claim in isolation.

This thesis introduces two conceptual contributions.
The first is the \textit{quantum-implicit programming tier}: a formal characterisation of a class of quantum programming languages whose surface model is agnostic of the underlying physical quantum substrate.
It establishes a new position on the quantum programming language abstraction spectrum, above existing high-level languages.
This position is defined by the absence of mandatory quantum-mechanical vocabulary from the programmer-facing surface and is operationalised across all four incorporated publications.
The second is \textit{quantum bias}: a named, formally operationalised measure of the degree to which a language's surface model requires the programmer to understand quantum-mechanical vocabulary to write correct programs.
Quantum bias is introduced and operationalised in Chapter~\ref{ch:pub_qpl_review}.

Four incorporated publications deliver the artefact contributions:
\begin{itemize}
    \item \textbf{Chapter~\ref{ch:pub_qpl_review} (quantum programming language systematic review):} the first systematic review of quantum programming languages using a novel abstraction-spectrum taxonomy and evaluation rubric.
    It establishes the state of the field, formally introduces and operationalises quantum bias and provides the evidential basis for the quantum-implicit tier as an unoccupied position in the design space.

    \item \textbf{Chapter~\ref{ch:pub_language_design} (language specification):} a formal specification of a quantum-implicit programming language, including its syntax, operational semantics and type system.
    It demonstrates that a quantum-implicit programming model can be precisely defined and that representative quantum algorithms can be expressed without quantum-mechanical vocabulary.

    \item \textbf{Chapter~\ref{ch:pub_formal_theory} (type system and formal theory):}formal proofs of type safety for the core language (progress and preservation) and semantic preservation for the compilation pipeline.
    It establishes the internal correctness guarantees of the artefact.

    \item \textbf{Chapter~\ref{ch:pub_compilation} (compiler):} a compiler from the quantum-implicit language to gate-based quantum hardware.
    It demonstrates end-to-end executability: programs written in the language can be compiled and run on current NISQ hardware.
\end{itemize}

This thesis establishes that the quantum-implicit tier is formally inhabitable and that at least one correct, compilable instantiation exists.
It does not claim that the specific language designed here is the only viable design, or that it is optimal.